\chapter*{General Conclusion}
\addcontentsline{toc}{chapter}{General Conclusion and Perspectives}
\markboth{General Conclusion and Perspectives}{General Conclusion and Perspectives}

This PFE addressed a question that is both scientific and practical: how can an electricity-consumption forecasting system be made reliable when initially impressive performance is affected by target leakage and unrealistic evaluation? The work began with the replication of a hybrid DEPM architecture and reproduced classification scores close to those reported in the starting article. A critical audit then showed that direct target information, the physical relationship between active power, voltage, and intensity, non-shifted rolling features, preprocessing before splitting, and random temporal separation could explain the apparent performance. A simple logistic-regression diagnostic reached 99.4\% under the suspect formulation, demonstrating that architectural complexity was not the source of the score.

The main scientific contribution was the reformulation of the task. Energy consumption became a continuous future trajectory; windows were causal; train, validation, and test regions were chronological; transformations were fitted on training observations; predictions were evaluated in physical units; and seasonal baselines were mandatory. The research compared recurrent, convolutional, patch-based, inverted-Transformer, decomposition, and global-model families while preserving the distinction between completed, incomplete, historical, and production experiments.

The ECL benchmark found iTransformer to have the lowest normalized MAE among the four completed architectures at all stored horizons. Checkpoint selection, however, was based on the different Low Carbon London household-transfer problem. Global TFT was selected for 24-hour and 168-hour horizons. The 24-hour model reaches cold-start macro MAE 0.1849 kWh per hourly interval, 26.48\% below seasonal naive, and beats the baseline for 99.0\% of households. The 168-hour model reaches 0.1973 kWh per hourly interval, 20.77\% below seasonal naive, and beats the baseline for 98.40\% of households. Reanalysis finds central-80\% interval coverage of 80.80\% and 78.35\%, respectively, under the original research preprocessing. Monthly experiments remain explicitly experimental because household-level macro $R^2$ is negative.

The engineering contribution is EnergyAI, a Dockerized Next.js and FastAPI platform with PostgreSQL persistence, background workers, authenticated ownership, bounded CSV ingestion, 24-hour and 168-hour forecast capabilities, quantile output, history, alerts, actions, analytics export, and readiness checks. Packaged checkpoints are governed by versioned manifests and SHA-256 hashes. Current quality evidence includes 100 passing backend tests on isolated PostgreSQL, successful frontend linting, type checking, and production compilation, and a Playwright matrix with one non-reproduced WebKit interruption. The exact frozen LCL cohort was also evaluated through the rolling 336-hour serving normalization without retraining: macro MAE is 0.181945 kWh at 24 hours and 0.194138 kWh at 168 hours. These supplementary values remain separate from the frozen research headlines. The core release is now publicly hosted on Azure with managed HTTPS, scale-to-zero web and API containers, managed PostgreSQL, a private container registry, Key Vault secrets, authenticated administrator access, and both forecast artifacts warmed. This supports a controlled public PFE demonstration. It does not close the geographically representative efficacy gap or the remaining operational gaps in email, registration, Google authentication, cloud workers, durable avatar storage, centralized logging, backup restoration, load, accessibility, disaster recovery, and penetration testing.

Several perspectives follow directly from the limitations. First, the models should be evaluated on Moroccan household or organizational data before local performance claims are made. Second, weather forecasts and static site attributes can be added through controlled ablations. Third, quantile calibration may benefit from conformal methods. Fourth, monthly forecasting requires improved household-level generalization. Fifth, finalists should be repeated over multiple seeds with deterministic sampling, mean, and standard deviation. Sixth, drift monitoring and a governed retraining workflow should be added without allowing unreviewed artifacts to replace packaged models. Finally, the public deployment should add a custom domain, application security headers, managed identity for registry access, durable avatar storage, centralized monitoring, cloud workers, provider checks, backup-and-restore evidence, dependency-gate remediation, load testing, and disaster-recovery procedures.

The most important result of this work is not a percentage. It is an evidence-traceable transition from a task that revealed its answer to a system that states what it knows, what it predicts, how it was validated, and where its limits remain. That transition is the foundation on which future EnergyAI improvements should be built.
